\documentclass[10pt,reqno]{amsart}

\usepackage[T1]{fontenc}
\usepackage[utf8]{inputenc}
\usepackage{newtxtext,newtxmath}
\usepackage{microtype}
\usepackage{mathtools}
\usepackage{booktabs}
\usepackage{enumitem}
\usepackage[nameinlink,noabbrev]{cleveref}
\usepackage[a4paper,margin=28mm,headheight=15pt]{geometry}

\setlength{\parskip}{2pt}
\setlist[enumerate]{leftmargin=7mm,itemsep=2pt,topsep=3pt}
\allowdisplaybreaks
\emergencystretch=3em

\newtheorem{theorem}{Theorem}[section]
\newtheorem{proposition}[theorem]{Proposition}
\newtheorem{lemma}[theorem]{Lemma}
\newtheorem{corollary}[theorem]{Corollary}
\theoremstyle{definition}
\newtheorem{definition}[theorem]{Definition}
\newtheorem{remark}[theorem]{Remark}

\newcommand{\Z}{\mathbb Z}
\newcommand{\R}{\mathbb R}
\newcommand{\PD}{\operatorname{PD}}
\newcommand{\tr}{\operatorname{tr}}

\title[A worked instance of Liu's Euler-class obstruction]
{A Fully Explicit Worked Instance of Liu's Euler-Class Obstruction}

\author{Argus Open Research Project}
\address{Argus Open Research Project}
\date{August 31, 2026}

\subjclass[2020]{Primary 57K32, 57K18; Secondary 57R30, 37E30}
\keywords{taut foliation, Euler class, mapping torus, pseudo-Anosov map, origami surface, Thurston norm}

\begin{document}

\begin{abstract}
Yi Liu's general nonrealizability theorem and its one-periodic-orbit
criterion are already established in arXiv:2409.14504.  This note records a
completely explicit worked instance \((M_f,w)\) of that criterion.  The
manifold is the mapping torus of a specified
affine pseudo-Anosov automorphism of the genus-three Eierlegende
Wollmilchsau origami.  In saturated free coordinates,
\[
 H_1(M_f;\Z)\cong \Z^3\oplus\Z/2\oplus\Z/6,
 \qquad \PD(w)=(-2,-2,2).
\]
The class \(w\) is even, integral, and lies in the dual Thurston unit ball,
but it is not the Euler class of any smooth cooriented taut foliation on
\(M_f\).  The nonrealizability follows from Liu's one-periodic-orbit
obstruction after an exact enumeration of all fixed trajectories.  All
finite algebra, homology, fixed-point, Alexander-polytope, and affine
noncollision assertions were independently replayed using exact arithmetic.
This document supplies no new obstruction theorem and no new disproof of the
Euler class one conjecture.  It is an exact computational application of
Liu's result, not a classification of the taut-foliation branch or of the
other five branches of Kirby Problem~3.15.
\end{abstract}

\maketitle

\section{Introduction and statement}

Let \(M\) be a closed oriented hyperbolic three-manifold with
\(b_1(M)>0\).  Kirby Problem~3.15 asks which real cohomology classes occur as
Euler classes of six kinds of structures, the first being taut foliations
\cite{BKR}.  Parity and the dual Thurston norm give familiar necessary
conditions.  They are not sufficient for the taut-foliation branch.

The purpose of this note is to record one exact, externally checkable pair.
It combines the quaternion origami of Herrlich--Schmithuesen
\cite{HerrlichSchmithuesen}, explicit affine-lift formulas of
Matheus--Yoccoz \cite{MatheusYoccoz}, the Alexander--Thurston comparison of
McMullen \cite{McMullen}, and the obstruction of Liu \cite{Liu}.  The
ingredients are established results; the construction below is a specified
application with all finite choices made explicit.

Let
\[
 Q_8=\{\pm1,\pm i,\pm j,\pm k\},
 \qquad i^2=j^2=k^2=ijk=-1.
\]
The Eierlegende Wollmilchsau \(W\) is the square-tiled surface with one unit
square for each \(g\in Q_8\), horizontal gluing \(g\mapsto gi\), and
vertical gluing \(g\mapsto gj\).  Let \(\widetilde S,\widetilde T\) denote
the standard affine lifts whose derivatives are
\[
 S=\begin{pmatrix}1&0\\1&1\end{pmatrix},
 \qquad
 T=\begin{pmatrix}1&1\\0&1\end{pmatrix}.
\]
Write \(L_g\) for left deck translation by \(g\).  Define
\begin{equation}\label{eq:f-definition}
 f=L_{-i}\circ\widetilde T\circ\widetilde S^2
       \circ\widetilde T\circ\widetilde S^2,
\end{equation}
where the rightmost map is applied first.  Equivalently, the affine
generators are applied in the order \(SSTSST\), followed by \(L_{-i}\).

Let
\[
 M_f=W\times[0,1]/(x,1)\sim(f(x),0).
\]
The exact homology calculation below supplies a saturated basis
\((\text{time},u,v)\) of the free quotient of \(H_1(M_f;\Z)\).  We record a
class in \(H^2(M_f;\Z)\) by the coordinates of its Poincar\'e dual in this
basis.

\begin{theorem}\label{thm:main}
The manifold \(M_f\) is closed, oriented, and hyperbolic, with
\[
 H_1(M_f;\Z)\cong\Z^3\oplus\Z/2\oplus\Z/6.
\]
There is an even integral class \(w\in H^2(M_f;\Z)\) with
\[
 \PD(w)=(-2,-2,2),
 \qquad x^*(w)\leq1,
\]
where \(x^*\) is the dual Thurston norm, such that \(w\) is not the Euler
class of any smooth cooriented taut foliation on \(M_f\).
\end{theorem}

\section{The origami monodromy}

The square decomposition of \(W\) has eight faces and sixteen edge classes.
The commutator of the horizontal and vertical gluing permutations is the
central element \(-1\), producing four vertex classes.  Hence
\[
 \chi(W)=8-16+4=-4,
\]
so \(W\) is a closed oriented surface of genus three.

\begin{lemma}\label{lem:derivative}
The derivative of the map in \eqref{eq:f-definition} is
\[
 A=Df=TS^2TS^2=
 \begin{pmatrix}11&4\\8&3\end{pmatrix}.
\]
It has characteristic polynomial \(X^2-14X+1\) and eigenvalues
\(7\pm4\sqrt3\).
\end{lemma}

\begin{proof}
This is direct matrix multiplication.  The left deck transformation has
derivative the identity, \(\det A=1\), and \(\tr A=14\).
\end{proof}

The derivative is hyperbolic.  Since \(f\) is an affine automorphism of the
translation surface \(W\), it is a pseudo-Anosov standard representative
with orientable invariant foliations.  Thurston hyperbolization for closed
pseudo-Anosov mapping tori therefore makes \(M_f\) hyperbolic.

\section{Integral homology and the selected class}

The cellular chain complex of the eight-square model has sixteen oriented
edge generators.  Applying the explicit affine lifts to those generators,
then imposing the square-boundary and vertex relations, gives the induced
integral action on \(H_1(W;\Z)\).  Adjoining the suspension generator and the
relations from \(f_*-I\) gives a mapping-torus presentation whose Smith
factors are
\begin{equation}\label{eq:smith}
 1,\ 1,\ 2,\ 6,\ 0,\ 0,\ 0.
\end{equation}
Consequently
\[
 H_1(M_f;\Z)\cong\Z^3\oplus\Z/2\oplus\Z/6.
\]

The exact row and column operations also give three primitive homomorphisms
to \(\Z\).  Their determinant on the free quotient is \(\pm1\), so the
resulting coordinates \((\text{time},u,v)\) are saturated.  Define
\begin{equation}\label{eq:classes}
 \PD(a)=(1,-1,1),
 \qquad \PD(e_f)=(-4,0,0),
 \qquad w=e_f+2a.
\end{equation}
Thus
\[
 \PD(w)=(-2,-2,2)=2(-1,-1,1).
\]
In particular, \(w\) is an even integral class, not merely a rational point
of \(H^2(M_f;\R)\).  The time coordinate gives
\begin{equation}\label{eq:fiber-evaluation}
 \langle a,[W]\rangle=1.
\end{equation}

\section{The dual Thurston condition}

Let \(P\subset\R^3\) be the Newton polytope of the multivariable Alexander
polynomial in the saturated free coordinates.  Its support has unique bottom
and top vertices in the fiber direction,
\[
 p_-=(0,2,2),
 \qquad p_+=(4,2,2).
\]
McMullen's equality on the fibered cone identifies the canonical Euler
functional with
\[
 \PD(e_f)=p_--p_+=(-4,0,0),
\]
in agreement with \eqref{eq:classes}.

The dual Alexander unit ball is the difference body \(P-P\) in the
normalization used here.  Two support points satisfy
\begin{equation}\label{eq:w-difference}
 (-2,-2,2)=(0,2,2)-(2,4,0),
\end{equation}
so \(\PD(w)\in P-P\).  An independent facet computation gives twelve
primitive supporting inequalities; direct substitution places \(w\) on the
facet \(-x-y=4\).  It also checks every one of the 137 distinct support
differences against all twelve inequalities.

Because \(b_1(M_f)=3\), McMullen's global comparison
\[
 \|\phi\|_A\leq\|\phi\|_{\mathrm{Th}}
\]
applies.  Duality reverses the inclusion, giving
\[
 B_A^*\subseteq B_{\mathrm{Th}}^*.
\]
Together with \eqref{eq:w-difference}, this proves
\begin{equation}\label{eq:dual-thurston}
 x^*(w)\leq1.
\end{equation}
No equality assertion for \(x^*(w)\) is needed.

\section{Complete exclusion of one-periodic trajectories}

The four singularities of the Abelian differential are permuted by \(f\) as
two transpositions.  Thus no singularity is fixed.  On the base torus,
\[
 |\det(A-I)|=12,
\]
so there are exactly twelve base fixed points.  Exhausting all eight
\(Q_8\)-lifts over each base point yields precisely twelve regular fixed
points of \(f\).  Their suspension homology classes are as follows.

\begin{table}[ht]
\centering
\caption{Complete one-periodic suspension multiset.}
\label{tab:periodic}
\begin{tabular}{cc}
\toprule
Free homology class & Multiplicity\\
\midrule
\((1,-1,0)\) & 1\\
\((1,-1,2)\) & 1\\
\((1,0,0)\) & 8\\
\((1,1,-2)\) & 1\\
\((1,1,0)\) & 1\\
\bottomrule
\end{tabular}
\end{table}

For completeness, the enumeration proceeds by choosing a rational based path
to each fixed lift, applying the explicit free-group automorphism induced by
\(f\), Schreier-rewriting the resulting closed word in the eight-square
cover, and reducing it in the saturated mapping-torus basis.  The sole fixed
point lying on a square edge is approached through a fixed rational interior
waypoint, avoiding an unrecorded choice of edge sector.

Every regular fixed point has local index \(-1\).  Their total index is
therefore \(-12\), which agrees with the Lefschetz number
\[
 L(f)=2-\tr(f_*|H_1(W;\Z))=-12.
\]
As a separate check, the first fiber layer of the exact multivariable
Alexander polynomial encodes exactly the five classes and coefficients in
\Cref{tab:periodic}.  The enumeration is therefore complete, rather than a
bounded orbit search.

By \eqref{eq:classes},
\[
 \PD(a)=(1,-1,1),
\]
which does not occur in \Cref{tab:periodic}.  Hence no one-periodic
trajectory represents \(\PD(a)\).

\section{Liu's obstruction and proof of the theorem}

We use the following consequence of Liu's one-periodic-orbit criterion
\cite[Lemma~3.2]{Liu}.  Let \(F\) be a closed oriented surface and let
\(g\colon F\to F\) be a pseudo-Anosov standard representative.  If an
integral class \(c\) on the mapping torus satisfies
\[
 \langle c,[F]\rangle=1
\]
and \(\PD(c)\) is not represented by a one-periodic trajectory, then the
class \(e_g+2c\) is not the real Euler class of any weakly symplectically
fillable contact structure.

\begin{proof}[Proof of \Cref{thm:main}]
The closedness, orientation, hyperbolicity, and homology statement follow
from \Cref{lem:derivative} and \eqref{eq:smith}.  The saturated coordinate
calculation proves that \(w\) is even and integral, and
\eqref{eq:dual-thurston} proves the required dual Thurston bound.

Apply Liu's criterion with \(F=W\), \(g=f\), and the class \(a\) in
\eqref{eq:classes}.  The fiber evaluation is one by
\eqref{eq:fiber-evaluation}, while \Cref{tab:periodic} proves that
\(\PD(a)\) has no one-periodic representative.  Therefore
\[
 w=e_f+2a
\]
is not the Euler class of any weakly symplectically fillable contact
structure.

The Eliashberg--Thurston approximation theorem, in the form recalled and used
by Liu, perturbs any smooth cooriented taut foliation to a weakly fillable
contact structure without changing the Euler class of the oriented plane
field \cite{EliashbergThurston,Liu}.  A smooth cooriented taut foliation with
Euler class \(w\) would therefore contradict the preceding obstruction.
\end{proof}

\section{Exact computational certification}

The construction involves substantial but finite algebra.  The accompanying
certificate was produced and replayed using integer and rational arithmetic;
no floating-point inequality is used in the proof of \Cref{thm:main}.  The
independent replay checks the following items.

\begin{enumerate}
\item The full multiplication table of \(Q_8\), the eight-square cellular
  boundary maps, and the canonical affine actions of \(\widetilde S\) and
  \(\widetilde T\).
\item The derivative word \(SSTSST\), the deck postcomposition \(-i\), the
  integral chain action, all Smith factors in \eqref{eq:smith}, and saturation
  of the three free coordinates.
\item All base fixed points, all \(Q_8\)-lifts, the four singularity images,
  every based Schreier word, and every homology class in
  \Cref{tab:periodic}.
\item The Alexander support, its first fiber layer, all 137 support
  differences, the complete rational facet system, and the membership of
  \(w\) in the difference body.
\item Every hypothesis used in the application of Liu's criterion.
\end{enumerate}

There is a close published affine formula in Matheus--Yoccoz
\cite[Remark~2.9]{MatheusYoccoz}.  It does not give the same pair.  Conjugating
the derivative of \(f\) by \(S^2\) gives
\[
 B=\begin{pmatrix}3&4\\8&11\end{pmatrix},
\]
whereas the printed element has derivative \(-B\).  Thus the traces are
\(14\) and \(-14\), respectively.  An exact simultaneous intertwiner system
for conjugacy and inverse conjugacy has rank eight and nullity zero.  Moreover,
the same-derivative factor in the published formula acts trivially on the
normal deck group \(Q_8\), while the present map induces the nontrivial
automorphism \(\operatorname{Ad}_i\).  The corresponding mapping tori have
first Betti numbers three and five.  These independent exact discriminators
exclude the nearest identified collision.

\section{Scope and literature status}

\Cref{thm:main} gives one explicit nonrealizable Euler class in the smooth
cooriented taut-foliation branch of Kirby Problem~3.15.  It does not classify
all Euler classes in that branch.  It makes no assertion about tight contact
structures, universally tight contact structures, pseudo-Anosov flows,
quasigeodesic flows, or circle actions.

Liu's theorem already establishes the general nonrealizability phenomenon,
and Lemma~3.2 of that paper is the core exclusion used here.  The quaternion
origami, its deck group, the canonical affine lifts, McMullen's norm
comparison, and the taut-foliation approximation theorem are also prior work.
The present document merely fixes one particular pair \((M_f,w)\) and records
the finite calculations needed to check that Liu's criterion applies.

Accordingly, this note is classified as a worked application and known-result
reconstruction.  It is not presented as an original theorem, a new
counterexample principle, a new disproof of the Euler class one conjecture, or
a solution of Kirby Problem~3.15(a).

\section{Independent-referee checklist}

An external referee can audit the result in the following order.

\begin{enumerate}
\item Verify the eight-square \(Q_8\) gluing and the genus-three Euler
  characteristic.
\item Multiply the six derivative generators and check
  \(Df=\left(\begin{smallmatrix}11&4\\8&3\end{smallmatrix}\right)\).
\item Recompute the mapping-torus Smith factors and confirm that the displayed
  free coordinates are saturated.
\item Check the twelve base fixed points, all lifted fixed points, and the
  complete homology multiset in \Cref{tab:periodic}.
\item Recompute the Alexander support and verify both
  \(\PD(e_f)=(-4,0,0)\) and \(\PD(w)\in P-P\).
\item Check the hypotheses of Liu's Lemma~3.2 and the passage from taut
  foliations to weakly fillable contact structures.
\item Treat the document as an explicit application of Liu's published
  obstruction, not as an original result or a solution of Problem~3.15(a).
\end{enumerate}

\begin{thebibliography}{99}

\bibitem{BKR}
R.~\.{I}. Baykur, R.~C. Kirby, and D. Ruberman,
\emph{K3: A New Problem List in Low-Dimensional Topology},
Mathematical Surveys and Monographs, vol.~295,
American Mathematical Society, Providence, RI, 2026.

\bibitem{EliashbergThurston}
Y. Eliashberg and W. Thurston,
\emph{Confoliations},
University Lecture Series, vol.~13,
American Mathematical Society, Providence, RI, 1998.

\bibitem{HerrlichSchmithuesen}
F. Herrlich and G. Schmith\"uesen,
An extraordinary origami curve,
\emph{Math. Nachr.} \textbf{281} (2008), no.~2, 219--237.

\bibitem{Liu}
Y. Liu,
On the Euler class one conjecture for fillable contact structures,
preprint, 2025, arXiv:2409.14504v2.

\bibitem{MatheusYoccoz}
C. Matheus and J.-C. Yoccoz,
The action of the affine diffeomorphisms on the relative homology group of
certain exceptionally symmetric origamis,
\emph{J. Mod. Dyn.} \textbf{4} (2010), no.~3, 453--486.

\bibitem{McMullen}
C. T. McMullen,
The Alexander polynomial of a 3-manifold and the Thurston norm on cohomology,
\emph{Ann. Sci. \`Ecole Norm. Sup. (4)} \textbf{35} (2002), no.~2, 153--171.

\end{thebibliography}

\end{document}
